%% DVXR — IEEE conference paper (full POW scope, honest write-up).
%% Methodology and Results prose is drafted from committed artifacts and adversarially
%% number-checked against outputs/*, presentation/data/*, and docs/*. Result tables are
%% auto-generated by scripts/build_paper_tables.py (dvxr.eval.paper); every number traces to a
%% committed scoreboard/CSV. Prose states nothing the tables and honesty audit do not support.
\documentclass[conference]{IEEEtran}
\usepackage{booktabs}
\usepackage{graphicx}
\usepackage{amsmath}
\usepackage{amssymb}
\usepackage[hidelinks]{hyperref}

\title{When Learned Cross-Modal Fusion Helps and When It Harms:\\
Subject-Held-Out Evidence from Small Clinical BCI\\
and Wearable Cohorts}
\author{\IEEEauthorblockN{Abhignya Jagathpally}
\IEEEauthorblockA{DVXR Lab, University of North Texas}}

\begin{document}
\maketitle

\begin{abstract}
Under subject-held-out evaluation on small clinical BCI and wearable cohorts, when does learned
cross-modal fusion help versus harm relative to the best single-modality specialist, and can a
reliability-gated late-fusion rule provide a do-no-harm floor? On six mental-health and clinical
tasks, learned CACMF fusion is outperformed by the strongest specialist (Holm $p{=}1.0$). Among
fusion methods, a Super-Learner-style do-no-harm (DNH) rule beats CACMF on four of six tasks, but
the inner-CV floor is not a held-out guarantee at $N\le 60$ (EEGMAT $-15.6\%$). The only
complementary-signal win is meal events added to CGM on the same CGMacros subjects (30-minute RMSE
$13.33\rightarrow 12.99$\,mg/dL). Specialist floors include WESAD AUROC $0.955$ ($n=8$) and a
Mumtaz LaBraM AUROC $0.961$ reported only as an identity-leakage-confounded upper bound. Every
head is Platt-calibrated and scored with decision-curve analysis. This is screening and decision
support, not a diagnostic device.
\end{abstract}

\begin{IEEEkeywords}
multimodal fusion, subject-held-out evaluation, EEG foundation models, continuous glucose
monitoring, do-no-harm late fusion, decision-curve analysis
\end{IEEEkeywords}

\section{Introduction}
Biosignal papers routinely report accuracies above 90\% for EEG depression detection or wearable
stress recognition. Three evaluation failures recur. First, many of those numbers come from
segment-level splits in which windows from the same person appear on both sides, so the score
overstates performance on a genuinely unseen subject. Second, scores are rarely calibrated, so a
model can rank well and still emit probabilities that are systematically wrong at a threshold.
Third, learned cross-modal fusion is often assumed to dominate the best single-modality specialist
once the streams are heterogeneous and the cohort is small. That last assumption is a claim, not a
default, and it is the claim this paper tests.

The research question is: under subject-held-out evaluation on small clinical BCI and wearable
cohorts, when does learned cross-modal fusion help versus harm relative to the best
single-modality specialist, and can a reliability-gated late-fusion rule provide a do-no-harm
floor? Scope is the parent DVXR pipeline on public, labeled cohorts (WESAD, PhysioNet Non-EEG,
DEAP, EEGMAT, Mumtaz, CGMacros). A separate cross-cohort glycemic FusionModel that pairs clamp
EEG onto CGMacros people is out of scope and is not a Goal~3 result.

Help and harm are defined on the same folds. Fusion \emph{helps} when fused error is strictly
below the strongest non-fused specialist, with Holm-controlled tests. Fusion \emph{harms} when
relative error reduction on $1-\mathrm{AUROC}$ (or RMSE for glucose) is negative. A
reliability-gated do-no-harm (DNH) rule may admit a blend only when an inner-CV advantage clears
a subject-grouped reliability gate; we report whether that floor holds on held-out subjects at
$N\le 60$. Nature attention fusion~\cite{kumar2026attention} is the published success condition
we do \emph{not} have: same-person, synchronized EEG+ECG \emph{stress}, not glucose (Fig.~\ref{fig:nature}).

Four contributions follow. (1)~A uniform causal ingestion path, frozen-encoder specialists, and
calibrated heads under subject-held-out cross-validation, with Platt scaling, conformal
intervals, and decision-curve analysis~\cite{vickers2006dca}. (2)~A systematic negative: learned
CACMF fusion is outperformed by the strongest specialist on all six mental-health and clinical
classification tasks (Fig.~\ref{fig:fusion}). (3)~Among fusion methods, DNH beats CACMF on 4 of 6
tasks, with a documented finite-sample caveat that the inner-CV floor is not a held-out guarantee.
(4)~The only complementary-signal win is meal events added to CGM history on the \emph{same}
CGMacros subjects (Fig.~\ref{fig:glucose-ablation}). Every number traces to a committed
scoreboard. The Mumtaz LaBraM AUROC $0.961$ is an identity-leakage-confounded upper
bound~\cite{lin2026identity}, not a validated biomarker. The system is screening and decision
support, not a diagnostic device.

\section{Related Work}
Subject-disjoint biosignal numbers sit far below leaky headlines. For depression from EEG,
MDD-SSTNet reports 65.1\% leave-one-subject-out (LOSO) accuracy on MODMA~\cite{mddsstnet}, and a
GoogleNet CNN drops to 73.3\% under external validation~\cite{metin2024mdd}, while a segment-level
pipeline reports 93.7\% with subject leakage~\cite{yan2022mdd}. Wearable stress repeats the gap:
an ensemble on merged corpora reaches roughly 85\% under LOSO~\cite{vos2023stress}; segment-level
WESAD numbers exceed 94\%~\cite{ghosh2022stress}. On PhysioNet mental arithmetic,
Khanam et al.\ study the subject-independent setting~\cite{khanam2023workload}, whereas a 4\,s
segment pipeline reports 97.4\%~\cite{yedukondalu2025workload}. Even a high
\emph{subject-disjoint} resting-state EEG score can be identity geometry rather than a
biomarker~\cite{lin2026identity,tai2026leaking}.

Learned multimodal fusion is licensed when complementary streams are jointly observed on the
same person. Kumar et al.~\cite{kumar2026attention} report 95.7\% three-class
rest/stress/amusement accuracy from synchronized EEG+ECG on WESAD+CASE (CASE has no EEG). That
paper does not forecast glucose. Figure~\ref{fig:nature} is the contrast, not a method we port:
same-person complementary biosignals versus the unmatched small-$N$ clinical setting of this
work. Clinical fusion under joint observation or modular missingness ---
MedFuse~\cite{medfuse}, MedPatch~\cite{medpatch}, HAIM~\cite{haim}, MultiModN~\cite{multimodn},
and DrFuse~\cite{yao2024drfuse} --- motivates \emph{testing} a learned mixer; those papers are
not glucose evidence here. Reconstructing a missing heterogeneous stream from another cohort is
unsafe~\cite{huang2020fusion,wu2024mlmm}. MEDFuse~\cite{phan2024medfuse} assumes joint
observation at train and test time. No open cohort co-registers EEG, CGM, and EHR on one subject.

Frozen specialists are probes, not universal winners. Independent benches find no single EEG
foundation-model champion under subject-disjoint evaluation~\cite{labram,xiong2025eegfmbench,lee2025capable}.
On short-horizon CGM, GlucoBench and GlucoFM-Bench require a persistence baseline; task-specific
models beat generic time-series foundation models when CGM history is present, and wearable-only
forecasts without CGM-in lose to persistence~\cite{gluformer,sergazinov2024glucobench,lu2026glucofmbench}.
Clinical text uses a frozen clinical BERT~\cite{clinicalbert}. Reliability-gated late fusion
follows Super Learner provenance~\cite{vanderlaan2007superlearner}.

The gap is a subject-held-out, small-$N$, specialist-versus-fusion scoreboard that is allowed to
be negative, plus a reliability-gated late-fusion floor whose finite-sample failure is measured
rather than assumed away.

\section{Method}
The objects under comparison are a per-task specialist encoder, a learned CACMF mixer, and a
reliability-gated DNH late-fusion rule, all scored on the same subject-held-out folds. Fusion is
defined only on co-occurring streams (Section~\ref{subsec:ingestion}). Encoders set the floor
(Section~\ref{subsec:encoders}); the two fusion operators are specified in
Section~\ref{subsec:fusion}. Calibration, uncertainty, and tests live in
Section~\ref{sec:experiments}. Figure~\ref{fig:architecture} is the serving path.
\subsection{Co-occurrence Constraint and Canonical Schema}
\label{subsec:ingestion}

Genuine learned fusion is valid only where modalities co-occur on the same subject at the same
time. In the open data used here that holds for wearable physiology with CGM and meal context
(CGMacros), EEG with peripheral physiology (DEAP), and chest/wrist physiology (WESAD). No open
cohort co-registers EEG, CGM, and EHR on one subject. The serving router therefore refuses
unmatched clamp-EEG concatenated onto CGMacros people; that join is not a parent-DVXR Goal~3
object.

Every source maps onto one long-format event table before any encoder runs. The required floor is
thirteen columns (\texttt{subject\_id}, \texttt{session\_id}, \texttt{timestamp\_utc},
\texttt{source\_system}, \texttt{device}, \texttt{modality}, \texttt{channel}, \texttt{value},
\texttt{unit}, \texttt{sampling\_rate\_hz}, \texttt{quality\_flag}, \texttt{label\_name},
\texttt{label\_value}); dataset-specific extras are preserved after that floor. Windowing is
strictly no-future-leakage: features at prediction time $t$ use history $\le t$, and the target
is held out of the feature set. Splits are patient-disjoint, so temporal causality is paired with
the absence of subject leakage. This is a no-future-information property, not an interventional
causal-effect claim.

Galea and EMOTIV sessions ingested on 2026-06-08 validate the schema and plug-and-play path only:
they are single-subject, unused as training or validation labels, and contribute no reported
metric.

\subsection{Specialist Encoders (the Floor)}
\label{subsec:encoders}

Each task's strongest non-fused specialist is the floor that fusion must beat. Encoders are
chosen under subject-held-out evidence, not segment-level headlines. Foundation models are
adopted only where they earn a place against a classical floor.

For scalp EEG we freeze LaBraM Base (5.8M parameters) and fit a linear probe on training folds
only~\cite{labram}. The contrast is a band-power feature set plus gradient boosting. LaBraM was
pretrained at 200\,Hz with 200-sample patches; the cohorts here are sampled at 64\,Hz, so the
probe is under-resourced on rate and montage. High resting-state AUROC under subject-disjoint
cross-validation can still be identity geometry~\cite{lin2026identity}; the depression number is
reported later as an upper bound, not as a biomarker.

For wearable and peripheral physiology we use spectral and HRV/EDA features with a tuned
gradient-boosting classifier (WESAD, PhysioNet Non-EEG). Heavier time-series foundation models
were evaluated and rejected for this channel: they do not reliably beat a task-specialized tree
under leave-one-subject-out evaluation.

For continuous glucose we use causal CGM history plus meal macros on the same CGMacros subjects.
NeuroGlycemicNet is the multimodal residual-over-persistence forecaster used in the meal
ablation: missingness-aware per-modality experts, a sign-constrained carb-response kernel, and
split-conformal intervals. It is not treated as a standalone architecture paper. Gradient boosting
on the same features remains the point-forecast ladder winner at 30 minutes; the network is
justified by calibrated intervals, abstention, and native multimodal residuals, not by winning
the 30-minute RMSE race.

Clinical notes use a frozen Bio\_ClinicalBERT encoder~\cite{clinicalbert} with a logistic head and
a TF-IDF floor under note-held-out grouped cross-validation. The notes are not patient-aligned
with the EEG, wearable, or CGM cohorts, so this pathway is a standalone capability and is not
fused. The large-language-model layer explains encoder and forecast outputs only
(\texttt{explanation.predicts == False}); it is never a predictor.

\subsection{Fusion Operators under Test}
\label{subsec:fusion}

Three predictors are compared on the same subject-held-out folds: the strongest non-fused
specialist, learned context-aware cross-modal fusion (CACMF), and reliability-gated do-no-harm
(DNH) late fusion. Figure~\ref{fig:architecture} is the serving path for those objects.

CACMF builds one vector-quantization codebook per supplied modality and combines the quantized
latents with a selectable operator (early concatenation, intermediate projection, late-weighted
heads, attention pooling, or a masked cross-modal transformer~\cite{medfuse,medpatch}). The
scientific object is the fused representation versus the strongest non-fused opponent, not a
five-strategy bake-off. A missing modality is assigned a learned absent token and masked out of
the gate; it is not zero-imputed. When no modality clears quality and freshness gates, fusion
abstains.

DNH (\texttt{dnh\_gated} in \texttt{src/dvxr/bench/gated\_fusion.py}) does not jointly train a
mixer. An inner subject-grouped cross-validation scores a candidate library of single-modality
heads, linear concatenation, a gradient-boosted concatenation, and, when available, a frozen
foundation-model probe. Non-negative Super-Learner-style weights~\cite{vanderlaan2007superlearner}
are shrunk toward the best candidate by an amount tied to inner-CV noise. A blend is admitted
only if its inner-CV advantage clears one subject-grouped bootstrap standard error; otherwise
the method falls back to the best candidate. The contribution is the finite-sample
characterization at $N\le 60$, not a new oracle inequality. The floor is an inner-CV guarantee;
whether it holds on held-out subjects is an empirical question for Results.


\section{Experiments}
\label{sec:experiments}

Every number in Results is produced under one contract: the scoring model has not seen its test
subject. Execution is offline, CPU-default, and deterministic (\texttt{seed=7} through folds,
initialization, and bootstrap resamples).

\subsection{Cohorts and Tasks}
\label{sec:cohorts}

Mental-health and BCI tasks use Mumtaz 2016 resting EEG for MDD versus control ($n=58$)
\cite{mumtaz2016}; WESAD chest/wrist physiology for stress ($n=8$) \cite{wesad}; PhysioNet
EEGMAT rest-versus-arithmetic with 19-channel EEG plus ECG at 64\,Hz ($n=20$) \cite{eegmat};
PhysioNet Non-EEG wearable stress \cite{noneeg}; and DEAP EEG plus peripheral physiology against
self-assessment-manikin anxiety and arousal labels \cite{deap}. DEAP is a documented
cross-subject chance ceiling, not a product screen. Metabolic forecasting uses CGMacros under
patient-disjoint splits at 30--120\,min horizons \cite{cgmacros}. MTSamples notes are scored as
a standalone language capability and are not fused. MIMIC-IV demo mortality (15 events in 252
encounters) is honesty-excluded as a validated claim \cite{mimic}. No cohort co-registers EEG,
CGM, and EHR on one subject.

\subsection{Subject-Held-Out Protocol}
\label{sec:cv}

The fusion scoreboard uses repeated grouped cross-validation, \texttt{repeats=5}$\times$\texttt{folds=5}
(25 folds), \texttt{seed=7}. Specialist headline tables may use $3\times5$ (15 folds) where the
committed screener manifest says so (for example
\texttt{outputs/product/screeners/mumtaz\_depression/manifest.json}); those protocols are not
blended into one number. The grouping key is the subject identifier, so every window from a person
stays on one side of each split.

Window-level AUROC pools held-out windows; subject-level AUROC aggregates each held-out person
to one score. On Mumtaz, a dedicated identity audit
(\texttt{outputs/\_r2/depression\_identity\_audit.md}) decodes subject identity at accuracy
$0.888$ against a $1/58$ chance rate, and each subject carries exactly one diagnosis. Under that
confound the depression AUROC cannot be read as a transferable biomarker. Opponent selection
shares the reported folds (single-level, not nested CV); the resulting optimism is stated, not
hidden.

\subsection{Metrics, Calibration, Utility, and Tests}
\label{sec:stats}

Classification uses AUROC and error $1-\mathrm{AUROC}$. Relative error reduction is
$\mathrm{RER}=(\mathrm{base\_err}-\mathrm{prop\_err})/\mathrm{base\_err}$ against the strongest
non-fused opponent on the same folds. A $\ge 50\%$ RER bar was pre-registered; it is a hurdle,
not an expectation. Glucose uses RMSE and MAE in mg/dL against persistence, plus
leave-one-modality-out ablations. Classification probabilities are Platt- or temperature-scaled
and scored by expected calibration error on out-of-fold scores. Glucose uncertainty is
split-conformal PI-95 coverage. Decision-curve net benefit~\cite{vickers2006dca} at threshold
$p_t$ is
\begin{equation}
\mathrm{NB}(p_t) = \frac{\mathrm{TP}}{n} - \frac{\mathrm{FP}}{n}\,\frac{p_t}{1-p_t}.
\end{equation}
A screen is credited as useful at $p_t$ only when a one-sided 95\% bootstrap lower bound on
$\mathrm{NB}(p_t)$ exceeds the better of treat-all and treat-none (Table~\ref{tab:utility}).
Head-to-head tests are paired one-sided Wilcoxon signed-rank tests on shared folds, Holm-corrected
across tasks, with 95\% bootstrap confidence intervals on RER and Cliff's $\delta$.


\section{Results}
\label{sec:results}
Four empirical moves answer the research question: the specialist floor, learned-CACMF harm, the
DNH finite-sample caveat, and the one same-subject help case. Every number traces to a committed
scoreboard. EHR notes and MIMIC mortality are standalone or honesty-excluded and are not Results.
\begin{table*}[t]
\centering
\footnotesize
\caption{DVXR Screen headline results under subject-held-out cross-validation. Window-level AUROC traces to the committed scoreboard; subject-level AUROC aggregates each held-out subject's windows (reported only for single-class-per-subject tasks). Research-grade screening, not diagnosis.}
\label{tab:headline}
\begin{tabular}{p{2.7cm}p{3.0cm}p{2.5cm}p{2.5cm}p{0.9cm}p{2.7cm}}
\toprule
Task & Encoder & AUROC (window) & AUROC (subject) & ECE & Protocol / N \\
\midrule
Depression screen (MDD vs healthy) from resting EEG & real LaBraM EEG foundation model (frozen linear-probe) & 0.961 [0.942, 0.976] & 0.986 [0.966, 0.999] & 0.030 & 3x5 subject-held-out CV, n=58 \\
Acute-stress screen from wearable physiology & band-power physiology features + tuned GBM & 0.955 [0.930, 0.978] & within-subject: n/a & 0.052 & 2x5 subject-held-out CV, n=8 \\
Cognitive-workload screen (rest vs task) & ECG autonomic (band-power); LaBraM improves the EEG-only path & 0.740 [0.710, 0.770] & within-subject: n/a & 0.025 & 3x5 subject-held-out CV, n=20 \\
Stress screen from peripheral physiology & peripheral physiology (band-power, concat) & 0.892 [0.860, 0.910] & within-subject: n/a & -- & subject-held-out CV \\
\bottomrule
\end{tabular}
\end{table*}

\subsection{The Specialist Floor Fusion Must Beat}
\label{subsec:mentalhealth}

Table~\ref{tab:headline} and Figure~\ref{fig:headline} are the per-task floors. Each number is the
strongest non-fused specialist under the protocol named in that table; those protocols are not
blended with the $5\times5$ fusion scoreboard.

On Mumtaz, a frozen LaBraM probe reaches window AUROC $0.961$ $[0.942,\,0.976]$ and subject AUROC
$0.986$ $[0.966,\,0.999]$ ($n=58$, $3\times5$)~\cite{mumtaz2016}. Identity is decodable at
accuracy $0.888$ against a $1/58$ chance rate, and each subject has one diagnosis
(\texttt{outputs/\_r2/depression\_identity\_audit.md}). The pair is an
identity-leakage-confounded upper bound~\cite{lin2026identity}, not a validated biomarker.

WESAD wearable GBM reaches $0.955$ $[0.930,\,0.978]$ ($n=8$; the headline table uses $2\times5$)
\cite{wesad}. Peripheral Non-EEG stress is $0.892$ $[0.860,\,0.910]$ \cite{noneeg}. On EEGMAT the
paper-wide ceiling is ECG-autonomic $0.740$ $[0.710,\,0.770]$ \cite{eegmat}. LaBraM raises the
EEG-only path from $0.636$ to $0.663$ and still sits below ECG. DEAP anxiety $0.53$ and arousal
$0.548$ are at chance under subject-held-out evaluation \cite{deap}; they are honest negatives,
excluded from product claims.

Table~\ref{tab:sota} is protocol-labeled context only. Placing $0.961$ beside MDD-SSTNet LOSO
accuracy $0.651$~\cite{mddsstnet} is a confound warning, not a cross-protocol win. Segment-level
$0.937$~\cite{yan2022mdd} is leaky context.

\begin{table*}[t]
\centering
\footnotesize
\caption{DVXR (subject-held-out AUROC) beside published results on the same/comparable cohorts, each labeled with its protocol and DOI. Cross-subject (LOSO / subject-independent) is the honest bar; segment-level numbers carry subject leakage and are shown for context only, never as a head-to-head win across mismatched protocols.}
\label{tab:sota}
\begin{tabular}{p{2.4cm}p{2.2cm}p{3.0cm}p{2.2cm}p{2.3cm}p{2.6cm}}
\toprule
Cohort & DVXR (AUROC) & Published & Score & Protocol & DOI \\
\midrule
mumtaz\_depression & win 0.961 / subj 0.986 & MDD-SSTNet (Sinc-CNN) & 0.651 accuracy & LOSO (cross-subject) & 10.1093/cercor/bhae505 \\
mumtaz\_depression & win 0.961 / subj 0.986 & GoogleNet CNN & 0.733 accuracy & external-validation & 10.1177/15500594241273181 \\
mumtaz\_depression & win 0.961 / subj 0.986 & EEGNet & 0.937 accuracy & within-subject/segment & 10.1515/bmt-2021-0232 \\
wesad\_stress & win 0.955 & GB+ANN ensemble & 0.850 accuracy & LOSO (cross-subject) & 10.1016/j.jbi.2023.104556 \\
wesad\_stress & win 0.955 & GAF + DNN & 0.948 accuracy & within-subject/segment & 10.3390/bios12121153 \\
wesad\_stress & win 0.955 & EDA + spectrogram ML & 0.964 accuracy & within-subject/segment & 10.1142/S0129065724500278 \\
eegmat\_workload & ECG 0.740 / EEG-only LaBraM 0.663 & EMD + SVM (subject-independent) & -- accuracy & subject-independent & 10.1371/journal.pone.0291576 \\
eegmat\_workload & ECG 0.740 / EEG-only LaBraM 0.663 & R-LMD + BAO + ensemble & 0.974 accuracy & within-subject/segment & 10.1038/s41598-024-84429-6 \\
\bottomrule
\end{tabular}
\end{table*}

\subsection{Learned CACMF Fusion Harms or Ties at Chance}
\label{subsec:ablation}

On the $5\times5$ fusion scoreboard, learned CACMF is outperformed by the strongest non-fused
opponent on all six mental-health and clinical classification tasks
(Fig.~\ref{fig:fusion}, Tables~\ref{tab:fusion}--\ref{tab:comparative}). Relative error reduction
on $1-\mathrm{AUROC}$ is $-19.9\%$ on peripheral stress (95\% CI $-28.6$ to $-13.5$),
$-185.9\%$ on WESAD ($-441.6$ to $-80.5$), $-0.6\%$ on DEAP anxiety ($-6.4$ to $5.4$),
$-1.2\%$ on DEAP arousal ($-7.0$ to $4.7$), $-40.4\%$ on EEGMAT ($-56.7$ to $-26.8$), and
$-148.2\%$ on Mumtaz ($-224.4$ to $-86.7$). Every Holm-adjusted Wilcoxon $p$ is $1.0$; the
pre-registered $\ge 50\%$ bar is \texttt{False} everywhere.

In native AUROC, the specialist wins or ties every classification row: Non-EEG $0.892$ vs.\ fused
$0.871$; WESAD $0.955$ vs.\ $0.871$; EEGMAT $0.740$ vs.\ $0.635$; Mumtaz $0.918$ (MOMENT/sota on
this board; headline LaBraM is the $0.961$ upper bound) vs.\ $0.795$. DEAP is a near-chance tie
($0.534$ vs.\ $0.531$; $0.548$ vs.\ $0.542$). Where one stream already carries the signal (ECG
for arithmetic; resting EEG for Mumtaz), joint training adds negative value. Where the signal is
weak (DEAP), fusion ties at chance.

\subsection{Reliability-Gated Late Fusion as a Do-No-Harm Attempt}
\label{subsec:dnh}

Among fusion methods, DNH beats CACMF \texttt{rep:fused} on four of six tasks (stress $+10.8\%$,
WESAD $+28.2\%$, EEGMAT $+17.7\%$, depression $+52.9\%$ RER; $5\times5$, same folds as
\texttt{BENCHMARK\_FINDINGS.md}). Versus the best single modality it wins three of six (stress
$+30.9\%$, WESAD $+25.3\%$, depression $+14.0\%$). The held-out floor fails on EEGMAT
($-15.6\%$) and both DEAP tasks ($-3.6\%$, $-4.5\%$). Inner-CV safety is not a held-out
guarantee at $N\le 60$: on EEGMAT the inner loop over-trusted a concat/GBM candidate that
generalized worse than single-ECG.

A 1-SE selection rule (prefer the simplest candidate within one inner-CV SE) closes the EEGMAT
gap and kills the WESAD win; it is a net negative, so strict argmin stays the default
(\texttt{outputs/\_dnh\_mh/dnh\_ablation\_1se.md}). Adding LaBraM to the DNH library recruits the
foundation model on Mumtaz (AUROC $0.910\rightarrow 0.961$) and is not fooled on EEGMAT, where
ECG remains the ceiling (\texttt{outputs/\_dnh\_labram/benchmark\_scoreboard.csv}). On every
mental-health task the strongest opponent remains non-fused.

\begin{table*}[t]
\centering
\footnotesize
\caption{Strongest specialist versus learned CACMF on the product-claim tasks. DNH beating CACMF on 4 of 6 is reported in the text, not in these cells.}
\label{tab:fusion}
\begin{tabular}{p{3.2cm}p{5.5cm}p{2.4cm}p{2.4cm}}
\toprule
Task & Best method (ours) & Learned CACMF & Delta AUROC \\
\midrule
mumtaz\_depression & LaBraM EEG FM = 0.961 & 0.795 & +0.166 \\
wesad\_stress & band-power + GBM = 0.955 & 0.871 & +0.084 \\
eegmat\_workload & ECG autonomic (task best) = 0.740 & 0.635 & +0.105 \\
stress & peripheral physiology = 0.892 & 0.871 & +0.021 \\
\bottomrule
\end{tabular}
\end{table*}

\begin{table*}[t]
\centering
\footnotesize
\caption{POW Goal 3 ablation: best single modality vs the integrated multimodal fusion, under repeated subject/patient-held-out grouped CV (Holm-corrected). Integrated fusion wins only on glucose forecasting; single-modality wins or ties on every mental-health task.}
\label{tab:comparative}
\begin{tabular}{llrrrl}
\toprule
Task & Metric & Best single modality & Integrated fusion & Delta & Verdict \\
\midrule
Stress (PhysioNet) & AUROC & 0.892 & 0.871 & -0.021 & single-modality wins \\
Stress (WESAD) & AUROC & 0.955 & 0.871 & -0.084 & single-modality wins \\
Anxiety (DEAP) & AUROC & 0.534 & 0.531 & -0.003 & approx. tie \\
Arousal (DEAP) & AUROC & 0.548 & 0.542 & -0.006 & single-modality wins \\
Cognitive workload (EEGMAT) & AUROC & 0.740 & 0.635 & -0.105 & single-modality wins \\
Depression (Mumtaz) & AUROC & 0.918 & 0.795 & -0.123 & single-modality wins \\
Glucose forecast (CGMacros) & RMSE@30 (mg/dL, lower better) & 13.330 & 12.990 & -0.340 & integrated wins \\
\bottomrule
\end{tabular}
\end{table*}

\subsection{Fusion Helps Only under Same-Subject Complementary Signals}
\label{sec:glucose}

The only complementary-signal win is meal events added to CGM history on the same CGMacros
subjects (Table~\ref{tab:glucose}, Fig.~\ref{fig:glucose}). At 30 minutes, persistence RMSE is
$17.40$\,mg/dL and gradient boosting is $12.48$. NeuroGlycemicNet is $12.99$. On this committed
ladder, GBM wins point RMSE at every horizon. A separate file
(\texttt{outputs/\_r2/deep\_tabular\_result.csv}) reports a redesigned deep net with later-horizon
edges; those cells are not Table~\ref{tab:glucose} and are not drawn in Fig.~\ref{fig:glucose}.

Leave-one-modality-out under the same patient-held-out split
(Table~\ref{tab:glucose\_ablation}, Fig.~\ref{fig:glucose-ablation}) isolates the second modality.
CGM plus meals is RMSE $12.99$; CGM only is $13.33$ ($+0.34$\,mg/dL when meals are removed).
That $-0.34$\,mg/dL is the only integrated win in the study. Without CGM, RMSE rises to $34.09$
and the model abstains on about $49.7\%$ of windows. Wearable-only forecasts that lack CGM-in
lose to persistence~\cite{lu2026glucofmbench}. Instability detection AUROC $0.976$ (hypoglycemia)
and $0.981$ (hyperglycemia) at 30 minutes is supporting evidence, not a second thesis.

\begin{table*}[t]
\centering
\caption{Glucose forecasting model ladder on CGMacros (patient-held-out), RMSE (mg/dL) by horizon. Gradient boosting is best point RMSE at every horizon on this ladder. Lower is better.}
\label{tab:glucose}
\begin{tabular}{lrrrr}
\toprule
Model & RMSE@30min & RMSE@60min & RMSE@90min & RMSE@120min \\
\midrule
persistence & 17.400 & 26.790 & 32.640 & 36.450 \\
linear\_ridge & 12.820 & 22.170 & 26.910 & 29.440 \\
decision\_tree & 14.900 & 23.970 & 28.270 & 30.300 \\
random\_forest & 13.240 & 22.450 & 26.910 & 29.290 \\
gradient\_boosting & 12.480 & 21.650 & 26.450 & 28.710 \\
mlp & 12.610 & 21.770 & 26.660 & 29.340 \\
neuroglycemic\_net & 12.990 & 22.180 & 26.900 & 29.060 \\
\bottomrule
\end{tabular}
\end{table*}

\begin{table*}[t]
\centering
\caption{Glucose missing-modality ablation at the 30-min horizon (CGMacros, patient-held-out). Adding meal events to CGM lowers RMSE 13.33 to 12.99; removing CGM collapses accuracy and the model abstains on about 50\% of windows rather than guess.}
\label{tab:glucose_ablation}
\begin{tabular}{lrrrr}
\toprule
Scenario & RMSE (mg/dL) & MAE (mg/dL) & Coverage & Abstention \\
\midrule
CGM + meal events (observed) & 12.987 & 8.702 & 1.000 & 0.000 \\
without CGM & 34.092 & 26.495 & 0.503 & 0.497 \\
without meal events & 13.326 & 8.876 & 1.000 & 0.000 \\
\bottomrule
\end{tabular}
\end{table*}

Decision-curve utility (Table~\ref{tab:utility}) attaches to the winning specialist, not to CACMF:
depression is useful on a 10--75\% subject-level band; WESAD on 5--70\%; EEGMAT only on 35--60\%.
\begin{table*}[t]
\centering
\footnotesize
\caption{Clinical utility by decision-curve analysis (Vickers and Elkin, 2006): the decision-threshold band over which the screener's net benefit exceeds both treat-all and treat-none, the peak added net benefit, and its one-sided 95\% bootstrap lower bound. The useful band is bootstrap-gated so a chance advantage reads as not useful.}
\label{tab:utility}
\begin{tabular}{lllll}
\toprule
Task & Level & Useful band & Peak net benefit & Bootstrap 95\% LB \\
\midrule
mumtaz\_depression & subject & 10--75\% & 0.441 & 0.312 \\
wesad\_stress & window & 5--70\% & 0.142 & 0.103 \\
eegmat\_workload & window & 35--60\% & 0.134 & 0.079 \\
\bottomrule
\end{tabular}
\end{table*}


% Publication figures from paper/scripts/render_paper_figures.py.
% Captions folded from paper/harness/phase4_viz/FIGURES.md (honesty-aligned).

\begin{figure*}[t]
\centering
\includegraphics[width=0.98\textwidth]{figures/fig1_architecture}
\caption{\textit{DVXR serving path: ingest, encode, route or fuse, calibrate, and explain.}
Schematic of the locked parent-DVXR path (not a performance plot). Each modality is
encoded by a specialist and either routed to the winning single-modality head or
offered to the fusion library. The language model narrates a frozen number and never
computes risk.
\textit{Note:} Nature \emph{Sci.\ Rep.} attention fusion is the same-person EEG+ECG
stress success case and is not a glucose claim.}
\label{fig:architecture}
\end{figure*}

\begin{figure*}[t]
\centering
\includegraphics[width=0.98\textwidth]{figures/fig2_fusion_vs_single}
\caption{\textit{Learned cross-modal fusion versus the best single-modality specialist
on six subject-held-out mental-health and BCI tasks.}
Source: \texttt{presentation/data/goal3\_ablation/comparative\_performance.csv}.
Panel A: AUROC. Panel B: signed $\Delta$ (fusion $-$ best single). Fusion loses or
ties on every row shown (Holm $p{=}1.0$). The depression bars are $0.918$ vs.\ $0.795$,
not the LaBraM $0.961$ identity-leakage-confounded upper bound.
\textit{Note:} The only complementary-signal win in that CSV is same-subject glucose
(CGMacros RMSE@30 $13.33\rightarrow 12.99$\,mg/dL), shown in Fig.~\ref{fig:glucose}.}
\label{fig:fusion}
\end{figure*}

\begin{figure*}[t]
\centering
\includegraphics[width=0.98\textwidth]{figures/fig3_glucose_ladder}
\caption{\textit{Patient-disjoint 30-minute glucose model ladder on CGMacros
(lower is better).}
Source: \texttt{presentation/data/goal2\_fusion/glucose\_model\_ladder.csv}
(\texttt{horizon\_minutes=30}). Gradient boosting (RMSE $12.48$\,mg/dL) edges
NeuroGlycemicNet ($12.99$). Persistence is the required baseline ($17.40$).
\textit{Note:} Reported as measured; the deep net is not dressed as a win. Meal-event
fusion is a separate same-subject ablation, not a depth claim.}
\label{fig:glucose}
\end{figure*}

\begin{figure*}[t]
\centering
\includegraphics[width=0.98\textwidth]{figures/fig4_nature_contrast}
\caption{\textit{When attention-weighted fusion is licensed, and when it is not.}
Kumar et al.\ report a same-person, synchronized EEG+ECG stress result
(\emph{Sci.\ Rep.} 16:15188, 2026). That inductive bias is not licensed for unmatched
EEG$\times$CGM$\times$EHR triples; this paper refuses that join.
\textit{Note:} Nature fusion is a stress-method contrast, not a glucose method.
Diabetes-attention-fusion glycemic fused macro-F1 is excluded (stale).}
\label{fig:nature}
\end{figure*}

\begin{figure*}[t]
\centering
\includegraphics[width=0.98\textwidth]{figures/fig5_headline_specialists}
\caption{\textit{Headline specialists under subject-held-out evaluation
(window-level AUROC with intervals from Table~\ref{tab:headline}).}
Depression $0.961$ is an identity-leakage-confounded upper bound, not a validated
biomarker (subject identity decodable at $88.8\%$, $52\times$ chance).
\textit{Note:} Research-grade screening, not a diagnosis.}
\label{fig:headline}
\end{figure*}

\begin{figure*}[t]
\centering
\includegraphics[width=0.98\textwidth]{figures/fig6_glucose_ablation}
\caption{\textit{CGMacros missing-modality ablation at 30 min (patient-held-out).}
Source: \texttt{presentation/data/goal3\_ablation/leave\_one\_modality\_out\_cgmacros.csv}.
\textit{Note:} Meal events added to CGM are the only complementary-signal win;
removing CGM collapses RMSE and forces abstention. Values are those in the CSV.}
\label{fig:glucose-ablation}
\end{figure*}


\section{Discussion}
\label{sec:discussion}

Kumar et al.~\cite{kumar2026attention} license attention-weighted fusion where EEG and ECG are
synchronized on the same person for \emph{stress} (Fig.~\ref{fig:nature}). That paper does not
forecast glucose. This study does not have that success condition: no open cohort co-registers
EEG, CGM, and EHR, learned CACMF loses on six of six mental-health tasks, and meals help only
because they sit on the same CGMacros subjects. Nature's convex mix is not a glucose method. An
unmatched clamp-EEG$\times$CGMacros join is refused.

Specialists win on these small-$N$ BCI and wearable tasks because a strong unimodal cue is
already present and joint training on $N\le 60$ dilutes it: WESAD GBM, EEGMAT ECG, Mumtaz
band-power or LaBraM. DNH helps when the inner-CV joint gain is real and redundancy is low; it
cannot manufacture a biomarker on DEAP or beat ECG on EEGMAT.

Table~\ref{tab:utility} is specialist utility, not fusion utility. Depression is useful on a
subject-level 10--75\% band; WESAD on 5--70\%; EEGMAT only on 35--60\%.

Limitations follow the evidence, not a product tour. Cohorts are small (WESAD $n=8$), so
intervals are wide. EEG is 64\,Hz against LaBraM's 200\,Hz pretraining. The depression AUROC
needs within-subject label variation before it can be a biomarker. The DNH floor diverges from
held-out subjects; a 1-SE rule is not a free lunch. Opponent selection shares the reported folds
(single-level CV optimism). Galea and EMOTIV sessions are schema validation only. EHR notes and
MIMIC mortality are standalone or honesty-excluded; they are not fused. The system is not a
medical device (\texttt{validated\_for\_clinical\_use=false}).

\section{Conclusion}
\label{sec:conclusion}

Learned cross-modal fusion harms relative to the best specialist on all six subject-held-out
mental-health and clinical tasks. Reliability-gated late fusion beats CACMF on four of six and is
the defensible fusion default, but it is not a held-out do-no-harm guarantee at $N\le 60$. Fusion
helps when complementary signals co-occur on the same person: meal events on CGMacros lower
30-minute RMSE by $0.34$\,mg/dL. Future work is larger co-registered cohorts, nested
cross-validation, identity-adversarial depression controls, and native-rate EEG---not a replay of
an unmatched glycemic mixer.

\section*{Data Availability}
All cohorts are public research datasets (Mumtaz 2016~\cite{mumtaz2016}, WESAD~\cite{wesad}, PhysioNet
mental-arithmetic~\cite{eegmat}, PhysioNet Non-EEG~\cite{noneeg}, DEAP~\cite{deap},
CGMacros~\cite{cgmacros}, MTSamples~\cite{mtsamples},
and the MIMIC-IV demo~\cite{mimic}); no restricted patient data are used. Result tables regenerate from
the released code via \texttt{scripts/build\_paper\_tables.py}.

\section*{Ethics}
Secondary analysis of de-identified public research cohorts. The system is offline with no PHI egress
and is not a medical device; a positive screen is a prompt to consult a qualified clinician, not a
conclusion.

\section*{Author Contributions}
A.\ Jagathpally: conceptualization, methodology, software, analysis, and writing.

\section*{Conflicts of Interest}
The author declares no competing interests.

\section*{Funding}
No external funding.

\section*{AI Usage Disclosure}
An AI coding assistant was used for software implementation and for drafting and copy-editing this
manuscript, including a multi-agent draft-and-adversarial-verification pipeline in which every stated
number was re-checked against a committed source file. All experimental design, results, and their
interpretation are the author's; every reported number is machine-verified against a committed
scoreboard or a persisted model artifact.

\bibliographystyle{IEEEtran}
\bibliography{references}
\end{document}
